%% kapa: exam=1 round=24 year=2013 period=1 form=A subject=economics subject_ko=경제원론 seq=19 number=59
%% source: Q-Net 감정평가사 기출문제 (cst00309) artlSeq=5073802 file=extracted/1008/5073802/2013년 제24회 감정평가사 1차시험 기출문제/1교시(A형).hwp
%% conversion: kordoc 4.24.51 -> tools/gen1.py (auto); review=proofread-sonnet; flags=none
\begin{question}{59}
소규모개방경제에서 수입 소비재 A에 관세를 부과할 때 이 시장에 나타날 경제적 효과에 관한 설명으로 옳은 것은? (단, 국내 수요곡선은 우하향, 국내 공급곡선은 우상향하며, A의 국제가격은 교역 이전의 국내가격보다 낮다.)
\begin{choices}
  \item 국내 소비자의 잉여는 증가한다.
  \item 국내 생산자의 잉여는 감소한다.
  \item 국내 소비는 감소한다.
  \item 국내 생산자의 생산량은 감소한다.
  \item 국내의 사회적 총잉여는 증가한다.
\end{choices}
\end{question}
\answer{3}
